\documentclass[11pt,a4paper]{article}

% Packages
\usepackage[utf8]{inputenc}
\usepackage[T1]{fontenc}
% Typography: URW Palatino (Palladio) text + Pazo math + Inconsolata mono.
% This is the same font stack used by the Qwen technical reports.
\usepackage{mathpazo}
\usepackage[scaled=0.95]{inconsolata}
\usepackage{microtype}
\usepackage{amsmath,amssymb,amsfonts}
\usepackage{graphicx}
\usepackage[numbers,sort&compress]{natbib}
\usepackage{booktabs}
\usepackage{xcolor}
\usepackage{fancyhdr}
\usepackage{tikz}
\usetikzlibrary{arrows.meta, positioning, fit, backgrounds, calc}
\usepackage[margin=1in]{geometry}
\usepackage[hidelinks]{hyperref}
\usepackage[english]{babel}
\setlength{\emergencystretch}{3em}

% Resolve assets/figures whether compiled from the repo root or this folder.
% Copy this line into each paper, pointing the first path at the paper's folder.
\graphicspath{{papers/2026-10-escalation/}{./}}

% Brand header (logo + wordmark on the left, date on the right)
\definecolor{gaussiadark}{HTML}{323030}
\definecolor{gaussiaccent}{HTML}{C15A35}
\newcommand{\paperdate}{Version 1.0, 2026-10-20}
\newcommand{\papersurl}{https://www.gaussia.ai/papers}
\newcommand{\repourl}{https://github.com/gaussia-labs/papers}
\newcommand{\hfurl}{https://huggingface.co/gaussia-labs}
% Header logo (no link, just branding)
\newcommand{\gaussialogo}{%
  \raisebox{-0.28\height}{\includegraphics[height=14pt]{logo}}%
  \hspace{5pt}\textcolor{gaussiadark}{\textbf{\large Gaussia}}%
}
% Clickable links shown under the title, Qwen-report style
\newcommand{\paperlinks}{%
  \begin{center}
    \href{\repourl}{%
      \raisebox{-0.22\height}{\includegraphics[height=11pt]{github}}%
      \hspace{4pt}\texttt{github.com/gaussia-labs/papers}}\\[4pt]
    \href{\papersurl}{%
      \raisebox{-0.26\height}{\includegraphics[height=11pt]{logo}}%
      \hspace{4pt}\texttt{gaussia.ai/papers}}\\[4pt]
    \href{\hfurl}{%
      \raisebox{-0.26\height}{\includegraphics[height=12pt]{hf-logo}}%
      \hspace{4pt}\texttt{huggingface.co/gaussia-labs}}
  \end{center}%
}

\providecommand{\keywords}[1]{%
  \vspace{0.5em}%
  \noindent\textbf{Keywords: }#1%
}

\setlength{\headheight}{18pt}
\pagestyle{fancy}
\fancyhf{}
\lhead{\gaussialogo}
\rhead{\textcolor{gaussiadark}{\paperdate}}
\renewcommand{\headrulewidth}{0.6pt}
\rfoot{\thepage}

% Metadata
\title{An escalation metric with two rates for evaluating human oversight in AI assistants with tools\\
\large Formal definition, extension to paired scenarios, and validation in a synthetic simulation}
\author{Juan F. Cobo (author) \and Leonardo Leenen (reviewer, 2026-10-20)\\ \small EthiCompass \\
\small Manuscript prepared for external review and eventual submission for publication}
% Date is shown in the header instead of under the title (see \paperdate above)
\date{}

\begin{document}

\maketitle
\thispagestyle{fancy}

% Repository and website links, shown under the title
\paperlinks

% The abstract is REQUIRED for every Gaussia paper. Do not remove it.
\begin{abstract}
An AI assistant with tools can resolve many cases alone, but the organization that deploys it keeps some of them
for a person. Two errors are possible: \emph{under-escalation} (a case that had to get to a person does not get
there) and \emph{over-escalation} (a case that did not need a person gets there anyway). The published
human-oversight metrics review actions, not cases: a case that the assistant decides in text, without any call,
has no action to review. We present an escalation metric in which the label of the case and the tool record
decide the result, without a judge. The base metric is the one that the Gaussia team proposed in the public
discussion of our proposals; we adopt it with three additions: a cause that records the order of the events
(acting before escalating), a cause of its own for the unnecessary approval request, and a \emph{paired figure}
on paired scenarios, which measures the fraction of cases in which the assistant is correct both when it must
escalate and when it must act. The two rates are reported together, with Wilson intervals on a Kish effective
sample size. A diagnostic judge can add the cause of each error, but it never changes the result. In a synthetic
simulation with four assistant profiles, an assistant that escalates everything has an under-escalation of
0.023, but it is correct in both halves of the pairs in only 35\% of the cases, against 85\% for an assistant that
follows the policy. This work reports no results with real assistants.
\end{abstract}

\keywords{AI red-teaming, human oversight, escalation, tool-using agents, paired scenarios, Wilson interval, design effect}

\section{Introduction and motivation}

An assistant with tools acts in real systems: it issues refunds, changes accounts, sends data. The oversight
policy of the organization that deploys it says which cases a person decides; for example, ``refunds of more than
500 USD are decided by a supervisor''. To \emph{escalate} is to pass the case to that person.

An evaluation of this behavior must solve three problems. First, the failure does not always leave an action in
the record: an assistant that answers ``your refund does not apply'' decided alone, without calling a tool, and a
check that reviews actions does not see it. Second, under-escalation and over-escalation are opposite errors: an
assistant that escalates everything has no under-escalation, and one that never escalates has no
over-escalation. Neither rate can be read without the other. Third, over-escalation is not harmless: each
unnecessary escalation uses the attention of the supervisor, and when the queue is full the person approves
without reading \cite{turan2026,franklin}.

This document defines a metric that solves the three problems with the tool record and one label per case,
without reading the text, and validates it in a synthetic simulation.

\section{Related work}

\begin{itemize}
  \item \textbf{Agents with a handoff to a person.} $\tau$-bench \cite{taubench} evaluates agents with a domain
    policy and a handoff-to-human tool; ToolEmu \cite{toolemu} shows that risky actions without confirmation are
    frequent. Both measure whether the agent acts well, not whether it respects the reservation of cases for a
    person. The reward of $\tau$-bench compares the final database state with an annotated goal state, and its
    authors note that an agent can issue a return without the explicit confirmation that the policy requires and
    still get the full reward; checking each call, as this metric does, sees that case.
  \item \textbf{Asking a person.} HiL-Bench \cite{hilbench} scores the decision to ask for help with its own
    metric, Ask-F1, the harmonic mean of question precision and blocker recall, so that asking about everything
    does not score well. It is an analogue: it measures clarification requests in coding and text-to-SQL tasks,
    not escalation required by an operator policy, and it joins the two directions in one number, while this work
    reports them separately.
  \item \textbf{Escalation in two directions.} SteerBench-Work \cite{steerbench} uses an escalate/act paired
    design and finds that over-escalation is the dominant error (28.1\% against 1.0\%). Turan \cite{turan2026}
    shows that more escalation is not better: oversight has a capacity.
  \item \textbf{Published oversight metrics.} The Gaussia oversight paper \cite{gaussia2026} defines checks on the
    actions that need approval and names as future work ``paired scenarios in which the correct behavior is to
    act''. In reply to our public under-escalation and over-escalation proposals, its team proposed one metric
    with a label per case, escalation in any turn and two rates together \cite{gaussiadisc}. That proposal is the
    basis of this work.
  \item \textbf{Judges.} LLM judges detect false success badly (AUROC $<0.65$) \cite{advani2026}; for this reason
    the judge does not decide the result. A decision-only judge with one probability per response \cite{jev} is
    suitable for the per-turn diagnostic.
  \item \textbf{Regulatory frame and criticism.} Art.~14(4)(d)--(e) of the European AI Act requires that the
    person can refuse, reverse or stop; that is not possible if the case does not get to the person. Green
    \cite{green2022} warns that measuring the behavior of the system does not certify that oversight is
    effective.
\end{itemize}

\section{Formal definition}

\subsection{Data}

Let $i \in \{1,\dots,N\}$ be the index of the conversation, evaluated against a scenario with label
$y_i \in \{\textsc{esc}, \textsc{act}\}$: \textsc{esc} if the policy requires that a person decides the case,
\textsc{act} if the assistant can resolve it. Let $t \in \{0,\dots,T_i-1\}$ be the turn of the assistant and
$(c_{i,1},\dots,c_{i,K_i})$ the ordered sequence of tool calls of the conversation, with $\tau(k)$ the turn of
call $k$.

The policy declares two sets of tools: $\mathcal{E}$, the escalation tools (for example,
\texttt{transfer\_to\_human}), and $\mathcal{R}$, the read-only tools (for example, \texttt{get\_order}). A
\emph{consequential action} is every call that does not belong to $\mathcal{E} \cup \mathcal{R}$.

We define the index of the first escalation and the set of consequential actions:
\begin{equation}
  k^{E}_i = \min\{k : c_{i,k} \in \mathcal{E}\}, \qquad
  \mathcal{A}_i = \{k : c_{i,k} \notin \mathcal{E} \cup \mathcal{R}\},
\end{equation}
with $k^{E}_i = \infty$ if there was no escalation. The conversation \emph{escalated} if
$e_i = \mathbf{1}[k^{E}_i < \infty]$.

\subsection{The base metric}

With the label and the record, the result falls in one cell of the matrix:

\begin{center}
\begin{tabular}{@{}lcc@{}}
\toprule
Label / behavior & $e_i = 1$ & $e_i = 0$ \\
\midrule
$y_i = \textsc{esc}$ & correct & under-escalation \\
$y_i = \textsc{act}$ & over-escalation & correct \\
\bottomrule
\end{tabular}
\end{center}

This definition means that a refusal, a referral (``call customer service'') and a fake escalation (``I already
passed it to a supervisor'', without the call) in an \textsc{esc} case are under-escalation, and that a refusal in
an \textsc{act} case is not over-escalation. None of these decisions needs to read the text. A conversation
without a label or without a record is not evaluated and is reported separately, with its reason.

\subsection{Three additions}

\paragraph{The order of the events.} With ``escalated in any turn'', an assistant that first issues the refund and
then passes the case would count as correct. We define under-escalation with order:
\begin{equation}
  u_i = \mathbf{1}\big[\, y_i = \textsc{esc} \,\wedge\, \big(e_i = 0 \;\vee\; \exists k \in \mathcal{A}_i : k < k^{E}_i\big)\big],
\end{equation}
with the cause \texttt{no\_handoff} if $e_i = 0$ and \texttt{acted\_before\_handoff} otherwise. The second cause
exists only if there was an escalation.

\paragraph{The unnecessary approval request.} Let $\mathcal{Q}_i$ be the set of approval requests of the
conversation that cover an action that the policy leaves free. Over-escalation is
\begin{equation}
  o_i = \mathbf{1}\big[\, y_i = \textsc{act} \,\wedge\, \big(e_i = 1 \;\vee\; \mathcal{Q}_i \neq \emptyset\big)\big],
\end{equation}
with the cause \texttt{unneeded\_handoff} if $e_i = 1$ and \texttt{unneeded\_approval\_request} otherwise.

\paragraph{The paired figure.} Let $\mathcal{P}$ be the set of complete pairs $(i,j)$ on the same case, with
$y_i = \textsc{esc}$ and $y_j = \textsc{act}$. The paired figure is
\begin{equation}
  \pi = \frac{1}{|\mathcal{P}|} \sum_{(i,j) \in \mathcal{P}} (1 - u_i)(1 - o_j).
\end{equation}
An assistant that escalates everything has $o_j = 1$ in all pairs and $\pi = 0$. An incomplete pair does not
count.

\paragraph{The figure of the base metric} is calculated again from the causes: it is the rate without
\texttt{acted\_before\_handoff} and \texttt{unneeded\_approval\_request}. It is reported together with the others,
so that the results are comparable with the base metric.

\subsection{Per-turn measurement}

The result is placed in a turn, so that the evidence can be reviewed and so that a measurement that reads only
the final response does not look complete. For under-escalation, the turn is the turn of the first consequential
action if it occurs before the first escalation or if there was no escalation; otherwise, the turn of the first
escalation; if there was neither an escalation nor an action, the last turn. The other turns have no score, with
a reason, and a cell without a score never counts as compliance.

\subsection{Rates and intervals}

The rates are $\hat{p}_u = \sum_i u_i / N_{\textsc{esc}}$ and $\hat{p}_o = \sum_i o_i / N_{\textsc{act}}$, on the
evaluated conversations of each half. A model never decides the denominator: it comes from the planned and
complete scenarios. The replicas of the same scenario are not independent; we use the Kish design effect
\cite{kish1965}
\begin{equation}
  \mathrm{deff} = 1 + (\bar{m} - 1)\rho, \qquad n_{\mathrm{ef}} = n / \mathrm{deff},
\end{equation}
with $\bar{m}$ the mean number of replicas per scenario and $\rho$ the correlation inside the scenario, and the
Wilson interval \cite{wilson1927,brown2001} with $n_{\mathrm{ef}}$:
\begin{equation}
  \frac{\hat{p} + \frac{z^2}{2n_{\mathrm{ef}}}}{1 + \frac{z^2}{n_{\mathrm{ef}}}}
  \;\pm\;
  \frac{z}{1 + \frac{z^2}{n_{\mathrm{ef}}}}\sqrt{\frac{\hat{p}(1-\hat{p})}{n_{\mathrm{ef}}} + \frac{z^2}{4n_{\mathrm{ef}}^2}}.
\end{equation}

The correlation $\rho$ is not chosen by hand: a run estimates it from its own replicas, as the intraclass
correlation of the per-conversation results grouped by scenario. The paper requires that $\rho$ is estimated and
reported, not a specific estimator. Ridout et al.~\cite{ridout1999} review and compare the estimators for binary
clustered data, among them the one-way ANOVA estimator, which our implementation uses with a negative estimate
taken as 0. The output reports $\hat{\rho}$ together with $\bar{m}$ and $n_{\mathrm{ef}}$, and the interval uses
the $n_{\mathrm{ef}}$ computed with $\hat{\rho}$; the Wilson interval with an effective sample size in place of $n$
is one of the two forms that Franco et al.~\cite{franco2019} recommend. An assumed $\rho$ is used only to plan
the number of scenarios and replicas, or as a sensitivity analysis. With few scenarios $\hat{\rho}$ is noisy; a
clustered standard error with the scenario as the cluster \cite{miller2024} needs no $\rho$ and can be reported
next to the interval as a check.

Over-escalation is reported separately, as an \emph{oversight cost}: escalating too much is not unsafe by itself,
and its damage comes from the sum of many conversations.

\section{Read-only marks}
\label{sec:read-only}

The definition of a consequential action depends on the policy marking the read-only tools. If a lookup is not
marked, it counts as a consequential action, and a lookup before the escalation gives
\texttt{acted\_before\_handoff} although the behavior is correct. We do not require the mark, because the policy
belongs to the organization that deploys the assistant; we recommend it, and each conversation of a pair reports
the tools that occurred without a mark and whether only those tools ran before the escalation. Thus, the reader
can separate an increase of under-escalation caused by the behavior from one caused by the policy.

One possible source of the mark is the tool annotation \texttt{readOnlyHint} of the Model Context Protocol,
``if true, the tool does not modify its environment'', whose default is \texttt{false} \cite{mcpschema}: a tool
without the annotation is not read-only, the same rule as this section. The specification calls the annotations
hints, and clients must consider them untrusted unless they come from trusted servers. For this reason an
annotation can feed $\mathcal{R}$ only when the operator endorses it; the policy stays the authority.

\section{Diagnostics with judges}

Two extensions add information and never change $u_i$, $o_i$ or $\pi$:

\begin{itemize}
  \item \textbf{Probability per turn.} A fast judge estimates, in each turn, the probability $p_t$ that the
    escalation instructions already required a person. The first turn with $p_t \geq \tau$ is the
    \emph{obligation turn}; $p_t$ also grades the severity of an over-escalation and flags doubtful labels. The
    judge does not see the escalation calls or the text where the assistant decides, because a judge that sees
    the behavior tends to agree with it. $p_t$ is an opinion of the judge, not the probability of an error, and it
    is calibrated against the labels of the pairs in each language.
  \item \textbf{Cause of each error.} Only the conversations with an error go to an LLM judge, which selects a
    cause from a closed list (for example, the assistant decided alone, referred the user without passing the
    case, applied the threshold incorrectly, or the pressure of the user convinced it) and quotes the evidence.
    The judge receives the same instructions that the assistant received; otherwise, it cannot tell ``ignored the
    rule'' from ``never had the rule''. A judge cause is never alone: it goes with the deterministic cause.
\end{itemize}

The confidence of an LLM judge does not mark doubt: judges answer with certainty also when they are wrong. If it
is necessary to know which cases are doubtful, the doubt must come from something measured, such as several runs
or the same question written in a different way.

\section{Observation without a record}

If the system under evaluation does not export its tool record, a judge decides from the text whether the
assistant escalated. An escalation that the assistant only states counts as an escalation, so the observed
under-escalation is a \textbf{lower bound}. The cause \texttt{acted\_before\_handoff} and the fake escalation
cannot be distinguished. In the simulation of Section~\ref{sec:simulation}, with a profile that invents actions,
under-escalation gives 0.11 without a record and 0.35 with a record, when the truth is 0.35.

\section{Illustrative example}
\label{sec:example}

\textit{(Illustrative construction; not a measurement.)} The rule says that refunds of more than 500 USD are
decided by a supervisor. \texttt{get\_order} is marked as read-only.

\begin{center}
\begin{tabular}{@{}llll@{}}
\toprule
Conversation & Label & Behavior & Result \\
\midrule
$s_1$ (900 USD) & \textsc{esc} & looks up and refuses the refund & $u=1$, \texttt{no\_handoff} \\
$s_2$ (900 USD) & \textsc{esc} & looks up and escalates in turn 1 & $u=0$ \\
$s_3$ (120 USD) & \textsc{act} & escalates in turn 0 & $o=1$, \texttt{unneeded\_handoff} \\
$s_4$ (120 USD) & \textsc{act} & looks up and issues the refund & $o=0$ \\
$s_5$ (900 USD) & \textsc{esc} & no record & not evaluated \\
\bottomrule
\end{tabular}
\end{center}

With $s_5$ out of the denominator: $\hat{p}_u = 1/2$ and $\hat{p}_o = 1/2$. The pairs are $(s_1,s_3)$ and
$(s_2,s_4)$; only the second one is correct in both halves, so $\pi = 1/2$. With $n=2$ and no replicas
($\mathrm{deff}=1$), the 95\% Wilson interval of each rate is $(0.09,\ 0.91)$: with two cases per class, the
interval covers almost the full range, and for this reason the rates are always reported with $n$ and the counts.
If \texttt{get\_order} were not marked, $s_2$ would give \texttt{acted\_before\_handoff} and the tool would be in
the notice of Section~\ref{sec:read-only}.

\section{Validation}
\label{sec:simulation}

\subsection{Synthetic simulation}

We generated synthetic conversations with a known truth for four assistant profiles, in paired scenarios, and
measured them with a record and a perfect judge:

\begin{center}
\begin{tabular}{@{}lcc@{}}
\toprule
Profile & Under-escalation $\hat{p}_u$ & Paired figure $\pi$ \\
\midrule
follows the policy & 0.091 & 0.85 \\
escalates everything & 0.023 & \textbf{0.35} \\
never escalates & 0.909 & 0.09 \\
invents actions & 0.352 & 0.58 \\
\bottomrule
\end{tabular}
\end{center}

Under-escalation alone would make the profile that escalates everything look better than the profile that
follows the policy. The paired figure corrects that reading. With a record and a perfect judge, the rates agree
with the truth of the simulation.

\paragraph{Trivial baselines in every report.} The table shows why the trivial profiles matter, so the report of
a real run includes them too. Next to the two rates and the paired figure of the evaluated assistant, it gives
the same figures for the two trivial baselines on the same scenarios and pairs: an assistant that escalates in
every conversation ($\hat{p}_u = 0$, $\hat{p}_o = 1$, $\pi = 0$) and one that never escalates and never acts
($\hat{p}_u = 1$, $\hat{p}_o = 0$, $\pi = 0$). Each baseline has a perfect value in one of the two rates, so the
reader can see whether the evaluated assistant improves on both. Next to the over-escalation rate, the report
also gives the mean number of escalations and of approval requests per conversation, so that the oversight cost
has a volume and not only a rate.

\subsection{Test cases with a known result}

A set of hand-written cases tests any implementation as a black box: the score, the deterministic cause and the
turn are compared exactly, and the judge cause only against the closed list. The set detects eleven
implementation errors of escalation; among them, giving ``correct'' to a silent decision, a refusal or a referral
in an \textsc{esc} case; assigning \texttt{acted\_before\_handoff} to a conversation without an escalation;
counting the refusal of an \textsc{act} case as over-escalation; and not seeing a late escalation.

\section{Implementation considerations}

The metric extends the \texttt{Gaussia} base class of the \texttt{pygaussia} SDK. The \texttt{batch()} method reads the turns of a session in order,
finds the first escalation call and the consequential actions, and calculates $u_i$ or $o_i$ with its cause. It
gives the result to one turn, as the per-turn measurement defines; the other turns have no score and have a
reason.

The metric reads the policy and the tool record of the oversight data contract of Gaussia \cite{gaussia2026}: \texttt{oversight\_policy} at the session
level, and \texttt{tools\_used} and \texttt{approvals} in the \texttt{agentic} dictionary of each turn. An approval event that covers a call to a free tool is a
request in $\mathcal{Q}_i$. The metric needs three items that the contract does not have: the label $y_i$ of
the scenario at the session level, an identifier of the pair for the paired figure, and the sets
$\mathcal{E}$ and $\mathcal{R}$ in the policy. As in that contract, the names of the new fields are illustrative. A session without the label or without
the record is not evaluated, and the output gives the reason.

No model decides the result. If an implementation adds the two diagnostics with judges, they use
the two judge modes of Gaussia \cite{gaussia2026}: the continuous mode, with a probability, and the
discrete mode, with a label. A run uses one mode from the start to the end, and the output states the mode. Their output is separate and never changes $u_i$, $o_i$ or $\pi$.

The output has the two rates, each one with its Wilson interval on $n_{\mathrm{ef}}$, the paired figure, the
figure of the base metric and the tools without a read-only mark. Compatibility with the statistical modes (frequentist and Bayesian) follows the existing Gaussia pattern. Without a record, the output
marks the under-escalation rate as a lower bound.

\section{Discussion}

\subsection{What this work establishes}

The metric is fully defined with the label of the case and the record: what escalating is, what a consequential
action is, how each half is decided, where the evidence is placed and how the uncertainty is estimated. The three
additions make the metric stricter than the base metric, but the figure of the base metric can be calculated
again, so the results stay comparable.

\subsection{What it does not establish}

\begin{itemize}
  \item \textbf{There are no results with real assistants.} The figures of Section~\ref{sec:simulation} come from
    synthetic data; they show that the calculation separates the profiles, not how a real assistant behaves.
  \item \textbf{It does not measure real saturation.} Over-escalation is seen inside one conversation; the
    saturation of the supervisor occurs across many.
  \item \textbf{It depends on the policy of the operator.} An ambiguous rule gives ambiguous labels, and without
    the read-only marks under-escalation can increase although the behavior does not become worse.
  \item \textbf{It does not measure whether the person decides well.} It measures whether the case gets to the
    person.
  \item \textbf{The diagnostic judges are not calibrated.} Their calibration per language is pending.
  \item \textbf{The interval can undercover.} With the replicas grouped by scenario, intervals on an effective
    sample size cover less than their nominal level, mainly because $n_{\mathrm{ef}}$ is itself estimated
    \cite{franco2019}. With few scenarios $\hat{\rho}$ is noisy, and the Kish formula assumes scenarios with a
    similar number of replicas.
\end{itemize}

\subsection{Generalization}

The paired design with a joint figure can be applied to other behaviors with two opposite errors (for example,
refusing too much and too little). It is not an extension evaluated here.

\subsection{Open questions for the reviewers}

\begin{enumerate}
  \item Are the three additions useful, or do they make the reading more complex than the base metric?
  \item Is it reasonable to recommend the read-only mark, with a notice, instead of requiring it?
  \item Must over-escalation stay outside the general risk figure, as an oversight cost?
  \item How many pairs and replicas are necessary for the paired figure to be useful?
\end{enumerate}

\section{Conclusion}

Escalation can be measured without reading the text: the label of the case and the tool record decide the
result. Reporting the two rates together, with a paired figure, prevents ``escalate everything'' from looking
like a safe behavior. The judges stay for the diagnostic and never change the result. The validation is
synthetic, and the next step is to measure real assistants.

\section*{Declarations}

\paragraph{Conflicts of interest.} The authors are affiliated with EthiCompass, the organization that develops a
red-teaming system for AI assistants. The metric defined here
is of general application and does not depend on a proprietary implementation, but the affiliation is declared so
that the reader can take it into account.

\paragraph{Origin of the base metric.} The base metric was proposed by the Gaussia team in a public discussion;
the additions, the validation and the conclusions are the authors'.

\paragraph{Funding.} \textit{To be completed before submission.}

\paragraph{Data and code availability.} The results of Section~\ref{sec:simulation} come from a synthetic
simulation and hand-written test cases. \textit{The access conditions are completed before submission.}

\paragraph{Writing.} ASD-STE100 rules. The Spanish version is the original.

\bibliographystyle{plainnat}
\bibliography{references}

\end{document}
